% Discussion. Numbers are quoted from the results sections; sources are given there.
\subsection{What the three data sets agree on}

The three samples measure the bias in different ways. The quote panel uses mid-quotes at fixed
leads. The Kalshi trade sample uses the prices at which orders executed. The Polymarket sample uses
fills by identified wallets. They agree on four points.
\begin{itemize}
  \item \emph{Longshots lose and favourites roughly break even.} On Kalshi, longshot sides return
  $-55\%$ per dollar at the mid. On Polymarket, longshot buys return $-54\%$ per dollar. In
  Kalshi's trade sample, taker longshot positions lose in 14 of 18 strata. Favourites earn between
  zero and a few per cent on both venues.
  \item \emph{The bias is small per contract.} In probability units the mispricing is a fraction
  of a cent to a few cents (\cref{sec:res-xs-limits}). Per dollar it is large only because the
  price is small. Sides quoted at 3--5\cents{} return $-54\%$ per dollar at the mid, so they win
  about half as often as their price implies, yet the absolute error is only about two cents.
  \item \emph{Small traders hold the longshots.} On both venues, longshots are a large share of
  small orders or small wallets and almost none of large orders or whales.
  \item \emph{Nobody is reliably on the other side.} Kalshi makers earn little more than takers
  overall, and on Polymarket no class of wallets is a significant net seller of longshots. The
  wallets that profit do not persist.
\end{itemize}
They differ in one respect. Polymarket's market-level calibration slope a day ahead, $0.99$
(SE $0.15$), is imprecise and cannot be distinguished from Kalshi's $1.13$, but the equal-weight
estimate, $0.97$ (SE $0.05$), is lower. Polymarket trades on a sub-cent grid near the extremes (89\% of
its fills below 4\cents{} or above 96\cents{} are at sub-cent prices)
% source: data/research/flb_lead/checks.json
and charged no fees for most of the sample, so both mechanical frictions of \cref{sec:framework} were
weaker there. The samples, price measures and periods also differ, so we do not read the gap as a
venue effect. On the 441 games observed on both venues, prices are almost identical.

\subsection{Interpreting the bias}

\paragraph{Mechanics set the band, not the position within it.} The one-cent tick and the fee
schedule explain why the bias cannot be traded away. A contract worth a tenth of a cent cannot be
quoted below a cent, and a taker who sells the longshot pays a spread and a fee that absorb the
whole premium. This is a limit to arbitrage in the sense of \citet{shleifervishny1997}, created
here by market design rather than by capital constraints. The band permits mispricing; it does not
require it. Without the tick-floor buckets the slope is unchanged, and at a one-hour horizon,
where the capital term of \cref{prop:band} vanishes, the bias is at least as strong as at thirty
days. Something keeps longshot prices near the upper edge of the band. That something is demand.
The structural analysis sharpens this reading (\cref{sec:res-structural}).
\begin{itemize}
  \item Market-weighted, the band covers the whole longshot gap in every group.
  \item Volume-weighted, the gap exceeds the band by 0.79\cents{} per contract ($t=2.7$). In the
  busiest markets, where spreads are tight, prices go beyond what frictions allow.
  \item On the few markets quoted on a 0.1\cents{} grid, 1--2\cents{} sides remain overpriced beyond
  their band.
\end{itemize}

\paragraph{Who supplies that demand.} The trade-level evidence identifies the buyers.
\begin{itemize}
  \item Longshots are bought disproportionately by small orders and small wallets, as
  \cref{hyp:pref} predicts.
  \item The buying is not concentrated in moments of retail attention. On Kalshi, volume surges
  bring \emph{fewer} longshots and fewer small orders, and only weekends and the last hour before
  expiry show a steeper bias.
  \item Small traders do not lose more per dollar than large ones on the longshots they buy. They
  bear most of the bias because they hold most longshots.
\end{itemize}
This is the pattern of lottery demand by individual investors in equity and option markets
\citep{kumar2009,bryzgalova2023}: the demand is for the payoff profile, not the result of worse
execution.

\paragraph{Information plays a small role.} Insider trading would make longshot flow informative
and longshot prices protective. Neither holds. No order-size class predicts outcomes on Kalshi,
small-order flow neither predicts nor reverts, and Polymarket's herding by small wallets moves
prices for an hour without predicting outcomes. In multi-outcome events Shin's insider share is
small, and the longshot field stays overpriced after the Shin adjustment.
The insider share implied by the quoted margin is mostly a restatement of the spread. The share
that best fits outcomes is no higher where private information is plausible, and it is lowest in
politics.

\paragraph{Preferences and beliefs.} Two explanations remain: preferences, either risk-love or
probability weighting, and heterogeneous beliefs. Prices and outcomes alone cannot separate them.
Both produce a monotone calibration curve, which a weighting function fits by construction
\citep{manski2006,gandhi2015}. Separating risk-love from weighting would need bet sizes or exotic
bets \citep{snowbergwolfers2010}. Three findings tilt the balance towards preferences.
\begin{itemize}
  \item \emph{One weighting function fits everywhere.} The fitted curvature is below one in every
  category (0.80 to 0.93), and a common value is not rejected under market weights. The pooled
  function predicts out of sample as well as group-specific ones. Its gain in log score over
  face-value prices is negligible, because most quotes lie in the middle of the price range, where
  prices are calibrated.
  \item \emph{Dispersion proxies fail.} The bias is not significantly steeper in events with more
  outcomes or more turnover. An external forecast does not reliably flatten it.
  \item \emph{The buyers are small.} Longshots are bought by small orders and small wallets, the
  traders that the evidence on lottery stocks and options associates with skewness preference
  \citep{kumar2009,boyer2014,bryzgalova2023}.
\end{itemize}
Belief dispersion among these traders is not ruled out. A minority of noise traders with dispersed
beliefs produces the same pattern \citep{gandhi2015}, and our proxies for dispersion are coarse.
What the evidence does rule out is a bias set by informed traders or by frictions alone.

\subsection{Implications}

\begin{itemize}
  \item \emph{Reading prices as probabilities.} Event-contract prices are informative but biased
  in the tails. Users who read them as forecasts should shrink longshot prices towards zero and
  favourite prices towards one, or apply the calibration function of \cref{sec:res-xs}. The
  correction matters most for multi-outcome fields and deadline contracts.
  \item \emph{Market design.} The tick floor alone accounts for about half of the loss in the
  cheapest bucket. A finer grid near the extremes would let prices of very unlikely outcomes fall
  towards their probabilities. The fee per dollar is largest on longshots, so fees shape who can
  trade against the bias but do not explain why it exists.
  \item \emph{Retail participation.} The losses on longshots fall mainly on small accounts, and
  wallet profits are concentrated in a few accounts that do not persist. This mirrors the
  distribution of outcomes in retail options trading \citep{bryzgalova2023}.
\end{itemize}

\subsection{Limitations}

\begin{itemize}
  \item \emph{Retail is proxied.} On Kalshi we observe order size, not identity, and makers are
  anonymous. Polymarket wallet classes use each wallet's history within the sampled markets only.
  \item \emph{Heavy tails.} Per-dollar longshot returns are dominated by rare wins. Pooled trade-level
  longshot means are imprecise and can change sign between samples, so we rely on per-stratum
  estimates, sign tests, calibration slopes and per-contract measures.
  \item \emph{Samples and periods.} The trade sample covers two and a half months of 2026. The
  quote panel applies the November 2026 fee schedule to its whole history, so net returns before
  2025 are indicative. The Polymarket sample excludes markets with more than \$6 million of
  volume, which carry 42\% of its frame's volume.
  \item \emph{Maker returns.} Returns to resting orders assume fills at the quote and are an upper
  bound, because resting orders are adversely selected.
  \item \emph{Multiple testing.} We report many cells. We therefore emphasise patterns that hold
  across categories, halves of the sample and data sets, rather than single cells.
  \item \emph{Polymarket US.} The data host of Polymarket's US exchange disallows our agent in
  its \texttt{robots.txt}, so we use only its published fee schedule.
\end{itemize}
